\section{Локализованный HPE. Разделение внутри пары}\label{sec:hpe}

Путь А ограничивает кривизну на рабочем шаре, но радиус этого шара
может быть сравним с расстоянием до решения. В HPE
(\emph{hybrid proximal extragradient}) мы делаем серию внешних
проксимальных шагов. Каждый из них решает более простую задачу
на малом шаре, который переносится вместе с текущим центром.
Внешняя схема использует вычисленную невязку подзадачи и потому
допускает неточное внутреннее решение~\cite{monteiro}.

В этом разделе алгоритму известны $L_{0,i},L_{1,i}$ и число $R>0$
с $\norm{x_0-z^\star}\le R$. Положим
\[
 L_{1,\max}=\max_i L_{1,i},\qquad
 r=\min\{R,1/L_{1,\max}\},\qquad
 \mathcal K=1+R^2/r^2,
\]
где $1/0=+\infty$. Рабочее множество на внешнем шаге из $x\in C$
равно $D=C\cap\ball_r(x)$, а регуляризованный оператор равен
\[
 F_\Gamma(z)=F(z)+\Gamma(z-x),\qquad \Gamma>0.
\]
У подзадачи на $D$ есть единственное решение $z_\Gamma^D$.
Существование следует из непрерывности и теоремы Брауэра для
отображения $z\mapsto \Pi_D(z-F_\Gamma(z))$ на компактном $D$.
Единственность следует из $\Gamma$-сильной монотонности.

\begin{figure}[htbp]
\centering
\begin{tikzpicture}[node distance=8mm]
  \node[slide box,text width=31mm] (center) {Внешний центр $x$\\
  $D=C\cap\ball_r(x)$};
  \node[slide box,text width=56mm,right=of center] (probe)
    {Проба с параметром $\Gamma$\\
     Проходы пути А на фиксированном $D$};
  \node[slide box,text width=33mm,right=of probe] (out)
    {Поправка $(u,v)$\\
    Новый центр $x^+$};
  \draw[slide arrow] (center) -- (probe);
  \draw[slide arrow] (probe) -- (out);
  \draw[slide arrow] ([xshift=14mm]probe.south) -- ++(0,-8mm)
    -| node[pos=.25,below,font=\small] {отказ: $\Gamma$ уменьшается}
    ([xshift=-14mm]probe.south);
\end{tikzpicture}
\caption{Один внешний шаг HPE. При отказе меняется регуляризация,
а центр и рабочее множество сохраняются.}\label{fig:hpe-composition}
\end{figure}

\subsection{Внутренний проход и сертификат невязки}

Обозначим $g_i=\nabla p_i$. По \eqref{eq:smooth} коэффициент
$L_{0,i}+L_{1,i}\norm{g_i(\cdot)}$ вдоль отрезка длины $t$
может вырасти не более чем в $\exp(L_{1,i}t)$ раз.
Действительно, градиент липшицев на компактном отрезке,
а разностное неравенство даёт для коэффициента верхнюю оценку
производной $L_{1,i}$, умноженную на текущий коэффициент, при
параметризации по длине. Интегрирование даёт указанную оценку.
Поэтому допустимы вычислимые константы
\[
 \widehat L_i(x)=e\bigl(L_{0,i}+L_{1,i}\norm{g_i(x)}\bigr),\qquad
 L_\Gamma=\Gamma+\sum_i(\widehat L_i(x)+M_i)
\]
для липшицевости $g_i$ и $F_\Gamma$ на $D$, соответственно.

В подзадаче заменяем каждую пару двумя уровнями $(p_i,0)$ и $(0,Q_i)$.
Добавку $\Gamma\norm{z-x}^2/2$ относим к одной гладкой функции.
Её градиент вычисляется напрямую. Применяем фиксированный проход
пути А с $\delta=\Gamma/24$. Для уровня с градиентом, липшицевым
с константой $L$, выбираем кривизну $\max\{\delta,2L\}$. На
операторном уровне берём $L=0$.
Этого достаточно, чтобы ни одна секущая не вызвала отказ.
По \eqref{eq:pa-energy} выход $y$ из начальной точки $a$ даёт
\[
 G_\Gamma(y,w)\le(\Gamma/8)\norm{a-w}^2,\qquad
 G_\Gamma(y,w)=G(y,w)+\frac{\Gamma}{2}
                    \bigl(\norm{y-x}^2-\norm{w-x}^2\bigr).
\]
В частности, $G_\Gamma(y,z_\Gamma^D)\ge\Gamma\norm{y-z_\Gamma^D}^2/2$
по выпуклости исходных функций и условию ВН для $F_\Gamma$.
Значит, каждый такой проход сокращает расстояние до $z_\Gamma^D$
вдвое. При старте из $x$ после
$m_\Gamma=\ceil{\log_2(128L_\Gamma/\Gamma)}$ проходов получаем
\begin{equation}\label{eq:hpe-inner}
 \norm{y-z_\Gamma^D}
 \le \eta_\Gamma:=\frac{\Gamma \norm{z_\Gamma^D-x}}{128L_\Gamma}
 \le \frac r{128}.
\end{equation}
Расстояние $\norm{z_\Gamma^D-x}$ входит только в анализ.
Сохранение этого множителя существенно, когда проксимальное
перемещение мало.

Для принятой пробы вычислим
\[
 u=\Pi_D\bigl(y-F_\Gamma(y)/L_\Gamma\bigr),\qquad
 n_D=L_\Gamma(y-u)-F_\Gamma(y),\qquad
 v=F(u)+n_D .
\]
По свойству проекции $n_D\in N_D(u)$, где
$N_D(u)=\{n:\ip{n}{w-u}\le0,\ w\in D\}$ --- нормальный конус.
Точка $z_\Gamma^D$ неподвижна для того же проекционного отображения.
Нерастяжимость проекции и липшицевость дают
\[
 \norm{u-z_\Gamma^D}\le2\eta_\Gamma,\qquad
 \norm{u-y}\le3\eta_\Gamma,\qquad
 \norm{F_\Gamma(u)+n_D}\le6L_\Gamma\eta_\Gamma .
\]
Если $\norm{u-x}<r$, то $n_D\in N_C(u)$. Для $w\in C$
достаточно подставить в нормальное неравенство
$u+t(w-u)\in D$ при малом $t>0$ и поделить на $t$.
Таким образом, $v\in F(u)+N_C(u)$.
Кроме того, $\norm{u-z_\Gamma^D}\le\norm{z_\Gamma^D-x}/64$,
так что $\norm{u-x}\ge63\norm{z_\Gamma^D-x}/64$.
Подстановка в оценку $\norm{F_\Gamma(u)+n_D}$ даёт
\begin{equation}\label{eq:hpe-res}
 \norm{u-x+\Gamma^{-1}v}\le \frac1{21}\norm{u-x}.
\end{equation}
При $z_\Gamma^D=x$ обе стороны равны нулю.

\subsection{Выбор регуляризации}

В начале внешнего шага берём
$\Gamma^{(0)}=\max\{\Gamma_{\min},8\norm{F(x)}/r\}$.
Если $F(x)=0$, можно вернуть $x$ как решение. На каждой пробе
вычисляем $y$ по \eqref{eq:hpe-inner}. Принимаем первую пробу
с $\norm{y-x}\ge r/4$ или $\Gamma=\Gamma_{\min}$.
При отклонении переходим к следующей пробе с
$\Gamma\gets\max\{\Gamma_{\min},\Gamma/2\}$.
Поправку $u$ вычисляем только для принятой пробы.

Для первой пробы ВН в $z_\Gamma^D$ с точкой сравнения $x$
и монотонность дают
$\Gamma\norm{z_\Gamma^D-x}\le\norm{F(x)}$, если перемещение
ненулевое. Значит, оно не больше $r/8$. При последующих пробах
важно, что $D$ остаётся тем же. Пусть $z_j$ --- точные решения
двух подзадач с $\Gamma_1\ge\Gamma_2>0$ и $d_j=\norm{z_j-x}$.
Если сложить их ВН и применить монотонность и неравенство
Коши--Буняковского, получится
\[
 0\le(d_2-d_1)(\Gamma_1d_1-\Gamma_2d_2).
\]
При $\Gamma_1>\Gamma_2$ это даёт
$d_1\le d_2\le(\Gamma_1/\Gamma_2)d_1$. При равных $\Gamma_j$
решения совпадают. Если предыдущая проба отклонена, её
точное перемещение меньше $r/4+r/128=33r/128$. Следующая
проба увеличивает его максимум вдвое. Поэтому на принятой
пробе $\norm{z_\Gamma^D-x}\le33r/64$, и обе точки $z_\Gamma^D,u$
строго внутри шара. Значит, точное решение регуляризованной
задачи является решением и на $C$, а поправка удовлетворяет
\eqref{eq:hpe-res} с нормалью к $C$. На пробе, принятой по длине,
дополнительно $\norm{u-x}\ge29r/128$.

На рис.~\ref{fig:hpe-radius} показан запас до границы для принятой
пробы. Алгоритм~\ref{alg:hpe} задаёт полный порядок вычислений.
Внутри он вызывает только процедуру \pass{} из алгоритма~\ref{alg:path-a}
на фиксированном $D$. Поиск радиуса пути А здесь не используется.

\begin{figure}[htbp]
\centering
\begin{tikzpicture}[x=12cm,y=1cm]
  \draw[thin] (0,0) -- (1,0);
  \foreach \pos in {0,.25,1}{
    \draw[thin] (\pos,-.09) -- (\pos,.09);
  }
  \node[below,font=\small] at (0,-.09) {$0$};
  \node[below,font=\small] at (.25,-.09) {$r/4$};
  \node[below,font=\small] at (1,-.09) {$r$};
  \draw[thick] (0,.53) -- ({33/64},.53);
  \draw[thin] ({33/64},.45) -- ({33/64},.61);
  \node[font=\small,above] at ({33/128},.53)
    {$\norm{z_\Gamma^D-x}\le33r/64$};
  \draw[thick] (0,1.08) -- ({17/32},1.08);
  \draw[thin] ({17/32},1.00) -- ({17/32},1.16);
  \node[font=\small,above] at ({17/64},1.08)
    {$\norm{u-x}\le17r/32$};
  \draw[dashed,thin] (1,.10) -- (1,1.12);
  \node[font=\small,align=center] at (1,1.39) {граница\\шара};
  \node[font=\small,below] at (.25,-.42) {порог для $\norm{y-x}$};
\end{tikzpicture}
\caption{Радиальные границы на принятой пробе. Отрезки показывают
верхние оценки расстояния до $x$.}\label{fig:hpe-radius}
\end{figure}

\begin{algorithm}[htbp]
\caption{Локализованный HPE с внутренним проходом пути А}
\label{alg:hpe}
\begin{algorithmic}[1]
\REQUIRE $x_0\in C$, $\rho>0$, distance bound $R$, smoothness parameters and $(M_i)$.
\STATE $x\gets x_0$, $S\gets0$, $W\gets0$,
  $\Gamma_{\min}\gets\rho$, $r\gets\min\{R,1/L_{1,\max}\}$.
\WHILE{$S<1/(2\rho)$}
  \STATE Compute $F(x)$ and $\widehat L_i(x)$. \textbf{If} $F(x)=0$, \textbf{return} $x$.
  \STATE $D\gets C\cap\ball_r(x)$,
     $\Gamma\gets\max\{\Gamma_{\min},8\norm{F(x)}/r\}$.
  \WHILE{true}
    \STATE Form $F_\Gamma$ and compute $L_\Gamma$,
       $m_\Gamma\gets\ceil{\log_2(128L_\Gamma/\Gamma)}$.
    \STATE Split the pairs into $2n$ levels and add the regularizer to one function.
    \STATE Run $m_\Gamma$ passes on fixed $D$, starting from $x$ and reusing each output as the next starting point. Use $\delta=\Gamma/24$ and the known local constants. Let $y$ be the final output.
    \IF{$\norm{y-x}\ge r/4$ or $\Gamma=\Gamma_{\min}$}
      \STATE \textbf{break}
    \ENDIF
    \STATE $\Gamma\gets\max\{\Gamma_{\min},\Gamma/2\}$.
  \ENDWHILE
  \STATE $u\gets\Pi_D(y-F_\Gamma(y)/L_\Gamma)$,
     $n_D\gets L_\Gamma(y-u)-F_\Gamma(y)$.
  \STATE $v\gets F(u)+n_D$, $\lambda\gets1/\Gamma$,
     $x\gets\Pi_C(x-\lambda v)$.
  \STATE $S\gets S+\lambda$, $W\gets W+\lambda u$.
\ENDWHILE
\RETURN $W/S$.
\end{algorithmic}
\end{algorithm}


\subsection{Глобальный зазор и число внешних шагов}

Положим $\lambda=1/\Gamma$, $e_H=u-x+\lambda v$ и
$x^+=\Pi_C(x-\lambda v)$. Из выпуклости, монотонности и
наличия нормали получаем $G(u,w)\le\ip{v}{u-w}$.
Раскрывая квадраты и пользуясь свойством проекции, находим
\begin{align}
 2\lambda G(u,w)
 &\le 2\lambda\ip{v}{u-w} \notag\\
 &\le \norm{x-w}^2-\norm{x^+-w}^2-\norm{u-x}^2+\norm{e_H}^2
 \qquad(w\in C).
 \label{eq:hpe-energy}
\end{align}
При $w=z^\star$ первый член слева неотрицателен.
Каждый длинный шаг уменьшает квадрат расстояния до $z^\star$
хотя бы на $(440/441)(29r/128)^2$. Для относительной
точности $\rho$ берём $\Gamma_{\min}=\rho$ и останавливаемся
при первом $S=\sum_k\lambda_k\ge 1/(2\rho)$. Если краткая
проба принята на нижнем пороге, её вес $1/\rho$ даёт остановку.
Так число принятых шагов составляет $\mathcal O(\mathcal K)$,
а $S\le3/(2\rho)$.

\begin{theorem}[Зазор локализованного HPE]\label{thm:hpe}
При приведённом критерии остановки взвешенный выход
$\bar u=S^{-1}\sum_k\lambda_k u_k$ удовлетворяет
\[
 G(\bar u,w)\le\frac{\norm{x_0-w}^2}{2S}
             \le \rho\norm{x_0-w}^2\qquad(w\in C).
\]
При $\rho=\varepsilon/(2R^2)$ это даёт абсолютную точность
$\varepsilon$ на $C\cap\ball_R(x_0)$.
\end{theorem}

Действительно, в \eqref{eq:hpe-energy} последние два члена в сумме
неположительны по \eqref{eq:hpe-res}. При суммировании промежуточные
квадраты сокращаются. Выпуклость $G(\,\cdot\,,w)$ обеспечивает
оценку для взвешенного среднего. Если уже найдено точное решение,
его зазор неположителен.

\subsection{Покомпонентные оракульные бюджеты}\label{sec:hpe-cost}

Пусть $A_i=L_{0,i}+L_{1,i}\norm{g_i(z^\star)}$.
Константы $A_i$ определяются решением и используются в анализе.
Знак $\OO$ ниже скрывает множители, логарифмические по параметрам
задачи и точности. Для абсолютного зазора $\varepsilon$ на шаре
сравнения получаем
\begin{align}
 N_{\nabla p_i}
 &=
 \OO\!\left(17^{2n-1}\left[
       \mathcal K(1+\sqrt{L_{1,i}R})
       +\sqrt{\frac{\mathcal K A_iR^2}{\varepsilon}}\right]\right),
       \label{eq:hpe-grad}\\
 N_{Q_i}
 &=
 \OO\!\left(17^{2n-1}\left[
       \mathcal K+\frac{M_iR^2}{\varepsilon}\right]\right).
       \label{eq:hpe-op}
\end{align}
Оценки включают отклонённые пробы. Для относительной точности
$\rho$ вместо $R^2/\varepsilon$ используется $1/\rho$.
Разделение касается ведущих членов точности и параметров конкретных
оракулов. Глубина $2n$, радиус $r$ и логарифмы зависят от всей задачи.

\subsubsection{Почему нет экспоненты по расстоянию}

На принятой пробе обозначим $z=z_\Gamma^D$.
Каждый внешний центр удовлетворяет $\norm{x-z^\star}\le R$.
Из регуляризованного ВН на $C$ следует $\norm{z-z^\star}\le R$:
складывая ВН в $z$ и $z^\star$, получаем
$\ip{x-z}{z-z^\star}\ge0$. Для остатка линейной модели
$D_{p_i}(z,z^\star)=p_i(z)-p_i(z^\star)
-\ip{g_i(z^\star)}{z-z^\star}$ имеем
\[
 0\le D_{p_i}(z,z^\star)\le
 \ip{F(z)-F(z^\star)}{z-z^\star}
 \le\Gamma\norm{z-x}\norm{z-z^\star}\le\Gamma rR.
\]
Первое верхнее неравенство получается суммированием
неотрицательных брегмановских остатков и вклада монотонных
операторов. Для второго используем оба ВН.

Положим $d_i=\norm{g_i(z)-g_i(z^\star)}$.
Функция $p_i(\cdot)-\ip{g_i(z^\star)}{\cdot}$ минимальна
в $z^\star$ на всём пространстве. При $d_i>0$ сделаем в анализе
один шаг по её градиенту из $z$ с длиной шага по времени
$1/[e(A_i+L_{1,i}d_i)]$. Длина отрезка, умноженная на
$L_{1,i}$, не больше единицы. Коэффициент квадратичной модели
на нём ограничен $e(A_i+L_{1,i}d_i)$, поэтому
\[
 \frac{d_i^2}{2e(A_i+L_{1,i}d_i)}\le D_{p_i}(z,z^\star).
\]
Здесь нужна гладкость на всём пространстве, поскольку пробная
точка анализа может покинуть $C$. Отсюда
$d_i\le2eL_{1,i}\Gamma rR+\sqrt{2eA_i\Gamma rR}$ и
$L_{1,i}d_i\le A_i+4eL_{1,i}\Gamma R$, используя $L_{1,i}r\le1$.
Вторая оценка получается разделением на случаи
$L_{1,i}d_i\le A_i$ и $L_{1,i}d_i>A_i$. При $d_i=0$ они
выполняются непосредственно. Теперь $\norm{z-x}\le r$,
поэтому рост коэффициента на отрезке от $z$ до $x$ даёт
\begin{equation}\label{eq:hpe-curv}
 \frac{\widehat L_i(x)}{\Gamma}
 \le e^2\bigl(2A_i/\Gamma+4e L_{1,i}R\bigr).
\end{equation}

\subsubsection{Сумма по пробам и шагам}

За один проход с параметром $\gamma$ уровни $(p_i,0)$ и $(0,Q_i)$
по \eqref{eq:pa-cost} дают разные масштабы
$1+\sqrt{\widehat L_i(x)/\gamma}$ и $1+M_i/\gamma$.
Оба умножаются на $\mathcal O(17^{2n-1}m_\gamma)$ за пробу.
Аналитический регуляризатор меняет градиентный масштаб лишь
на константу. Для любой предыдущей пробной $\gamma\ge\Gamma$
имеем $m_\gamma\le m_\Gamma$. Геометрическое уменьшение $\gamma$
даёт $\sum_\gamma \gamma^{-a}=\mathcal O(\Gamma^{-a})$
для $a=1/2,1$, включая возможное последнее усечение.

Число проб за шаг имеет порядок
$1+\log(2+\norm{F(x)}/(r\Gamma))$. Для него и для $m_\Gamma$
можно брать логарифмы общей величины
\[
 \mathcal L=2+nR/r+\sum_i A_i/\Gamma_{\min}
       +\frac{\norm{F(z^\star)}+R\sum_i M_i}{r\Gamma_{\min}} .
\]
Это следует из \eqref{eq:hpe-curv} и
$\norm{F(x)}\le\norm{F(z^\star)}+
 \sum_i(d_i+r\widehat L_i(x)+M_iR)$.
Базовый член на шаг несёт множитель $\log^2\mathcal L$,
члены с $\widehat L_i$ и $M_i$ несут $\log\mathcal L$.
Один запрос $F(x)$ и до двух запросов для поправки
добавляют постоянную работу на шаг.
Наконец, $\sum_k1/\Gamma_k=S$ и
$\sum_k\Gamma_k^{-1/2}\le\sqrt{KS}$, где $K$ --- число
принятых шагов. Подстановка $K=\mathcal O(\mathcal K)$,
$S=\mathcal O(R^2/\varepsilon)$ и \eqref{eq:hpe-curv}
даёт \eqref{eq:hpe-grad}--\eqref{eq:hpe-op}.
